\SourceLeaf{224}
\fancyhead[L]{\small GENERAL INVESTIGATIONS OF CURVED SURFACES}
\renewcommand{\Article}[1]{\par\medskip\Needspace{4\baselineskip}\pdfbookmark[1]{Art. #1}{superficies-art-#1}\centerline{\textbf{#1.}}\nobreak\smallskip}
\begin{center}
GENERAL INVESTIGATIONS\par\medskip
{\large OF CURVED SURFACES.}\par\medskip
\rule{.65in}{.3pt}
\end{center}
\Article{1}
\indent Investigations into the directions of various straight lines in space usually gain considerably in clarity and simplicity if we use an auxiliary spherical surface of radius \(=1\) drawn about an arbitrary centre. Each point on it is taken to represent the direction of the straight lines parallel to the radius that ends at that point. Since the position of every point in space is determined by three coordinates, namely its distances from three fixed planes perpendicular to one another, we must first consider the directions of the axes perpendicular to those planes. We shall denote the points on the spherical surface that represent these directions by \((1),(2),(3)\). The distance between any two of these points will therefore be a quadrant. We shall also assume that the axes are directed toward the sides on which the corresponding coordinates increase.

\Article{2}
\indent It may be useful to set out here a few propositions that are often used in questions of this kind.

I. The angle between two intersecting straight lines is measured by the arc between the points on the spherical surface that correspond to their directions.

II. The orientation of any plane can be represented by a great circle on the spherical surface whose plane is parallel to the given plane.
